% SPDX-License-Identifier: CC-BY-4.0
% Copyright (c) 2026 Martin Schröder <info@swedishembedded.com>

\section{Recursive self-improvement on a sealed exam}\label{sec:rsi}

The 27B generations of \cref{sec:generations} were promoted twice and did not beat their base on
one measurement. This section reports the four registered series that followed, RSI1 to RSI4, on one
question: \emph{does a model that teaches itself become better at every one of several consecutive
generations, on items it never saw, with no stronger model, teacher or human answer anywhere in its
training?} Each registration was committed before its first run, fixed its exam, criterion and
analysis, and records its outcome whatever it was.\cl{RS01,RS02} Three series did not meet their
criterion; RSI4 did, in all four of its lineages (\cref{tab:rsi,fig:rsi-items}).\cl{RS29}

\paragraph{What ``recursive self-improvement'' means here.} A series is self-taught when every
generation's training material is generated by the policy itself, the champion the previous generation
produced, and the only judge of that material's correctness is an executable oracle, here the Lua
interpreter: the policy writes the curriculum, never the labels, and no stronger model and no human
answer enters training. It improves when each generation beats the one it was trained from on an exam
it never saw. That is what is measured. It is \emph{not} measured whether the improvement compounds
because an improved model is a better improver, that is, whether generation $g$'s proposals teach more
than the same amount of interpreter-verified material proposed by the base or by a fixed earlier
policy; no series ran that control, and \cref{sec:documentation} found on a related task that who
writes the tasks did not change the gain.\cl{SL23}

\subsection{What every series shares}\label{sec:rsi-design}

\paragraph{The learner.} The task family is predicting what a Lua~5.1 program prints. The learner is
the loop of \cref{sec:documentation} reduced to this one kind: the policy, Qwen3-4B with thinking off,
writes programs about sections of the Lua~5.1 manual, the pinned Lua~5.1.5 interpreter runs each
program and its output becomes the reference, and the policy is trained on its own programs with the
interpreter's output as the target. From RSI2 on every model call and every training run goes through
Brain's API: one \code{brain serve} per lineage serves the proposals, the predictions and the dev and
retention exams, and trains each candidate as a Brain training job.\cl{RS09,RS10} A lineage is a chain
of generations from the base that differs from its siblings only in its seed; each series runs four
lineages of three generations. Inside the loop a candidate is accepted by a dev gate on 400 units
(100 generator programs and 300 programs the base proposed, held out from training) with a retention
set of 300 units; neither holds an exam item.\cl{RS09}

\paragraph{The exam.} The exam is 889 Lua~5.1 programs written by code from 25 standard-library
templates over randomised and edge-case arguments, the items the base never answered in four sampled
attempts under a neutral request; the reference of each is what the interpreter prints (twice,
equal).\cl{RS03} Because it was selected as what the base could not do, a count of items right on it
measures progress on that hard subset, not a model's accuracy on Lua programs in general. It was
selected once, in RSI1, and every series uses the same file (its SHA-256 begins
\code{60ac3ae44651} in every verdict).\cl{RS20,RS31} No model writes or answers an exam item to build
it, and no run of the learner reads it.

\paragraph{The examiner.} From RSI2 on the exam is held by an independent examiner, a separate program
operated beside the learner, and the learner's only contact with it is a one-way hand-off
(\cref{fig:examiner}).\cl{RS04} For each lineage and generation the learner writes a directory with
the list of every program the candidate's chain trained on and, last, a candidate record (base model,
base digest, adapter path and digest, serving flags, commits). The examiner then
\begin{enumerate}
\item serves the base and every candidate itself, on one pinned serving build (Brain
  \code{d11e91340}) and, in RSI4, at one KV-cache tier (f32), and marks a verdict invalid when the
  served base is not the declared checkpoint (identified from the checkpoint's own files) or the
  adapter does not have the declared digest;\cl{RS06,RS32}
\item excludes from each step's comparison every exam item that equals, exactly or after removing
  comments and whitespace, any program either compared model was trained on, and reports how many
  it excluded;\cl{RS05}
\item asks each item once under the direct-completion protocol: the program and its question as the
  only user message, no system prompt, one greedy completion with thinking off and, in RSI4, at most
  2048 tokens; the answer is the last fenced block of the reply, or the whole reply if it has none,
  compared with the reference token by token per line with whitespace runs collapsed and blank lines
  ignored; a reply the token limit cut is graded as written;\cl{RS32,RS36}
\item pairs the candidate with its predecessor item by item (a win: the candidate is right and the
  predecessor wrong; a loss: the reverse), computes the exact one-sided sign test over the discordant
  items, and decides the registered rule: a step holds when wins exceed losses at $p<0.05$, a lineage
  succeeds when all three of its steps hold, and the series meets its criterion when at least three of
  the four lineages succeed;\cl{RS02}
\item returns, per step, a verdict with aggregates only (items right, wins, losses, $p$, whether the
  step holds, exclusion and overlap counts, integrity flags) and, over the series, a plan verdict;
  each is signed with the examiner's HMAC-SHA256 key, so the learner, which does not hold the key,
  cannot alter or forge one. No item text and no item-level grade leaves the examiner.\cl{RS04,RS31}
\end{enumerate}
The items themselves are not secret: the exam file was released with RSI1's records, so the
separation rests on the learner's code never opening it and on the exclusion of overlapping programs,
not on secrecy.\cl{RS03}

\begin{figure}[tbp]
\centering
\begin{tikzpicture}[fig]
\tikzset{cell/.style={text width=29mm}}
% learner column
\node[proc=okSky!22, cell] (gen) at (0,0) {generation $g$:\\propose, label,\\train via Brain};
\node[verif, cell] (lgate) at (0,-22mm) {learner's gate:\\dev 400 units,\\retention 300};
% hand-off, between the two
\node[data, cell] (hand) at (37mm,-22mm) {\textbf{hand-off}:\\trained programs,\\then candidate};
% examiner, down the first column and up the second
\node[data, cell] (exam) at (76mm,0) {\textbf{sealed exam}:\\889 items the base\\never answered};
\node[verif, cell] (excl) at (76mm,-22mm) {drop each item\\equal to a trained\\program (also with\\comments removed)};
\node[proc=okOrange!20, cell] (serve) at (76mm,-46mm) {serve all models:\\one pinned build\\and KV tier,\\digests checked};
\node[proc=okOrange!20, cell] (grade) at (113mm,-46mm) {direct completion:\\greedy decoding,\\thinking off,\\last fenced block};
\node[decide, aspect=2, inner sep=0.5pt, text width=14mm] (rule) at (113mm,-22mm) {$W>L$,\\$p<0.05$?};
\node[data, cell] (verd) at (113mm,0) {\textbf{signed verdict}:\\counts, $p$, holds;\\no item, no grade};
\begin{scope}[on background layer]
\node[bound, fit=(gen)(lgate), inner sep=2.5mm] (L) {};
\node[bound, fit=(exam)(serve)(grade)(verd)(rule), inner sep=2.5mm] (E) {};
\end{scope}
\node[boundlabel] at (L.north west) {Learner};
\node[boundlabel] at (E.north west) {Independent examiner};
\draw[flow] (gen) -- (lgate);
\draw[flow] (lgate) -- (hand);
\draw[flow] (hand) -- (excl);
\draw[flow] (hand.south) |- (serve.west);
\draw[flow] (exam) -- (excl);
\draw[flow] (excl) -- (serve);
\draw[flow] (serve) -- (grade);
\draw[flow] (grade) -- node[lab, left, align=right] {wins,\\losses} (rule);
\draw[flow] (rule) -- node[lab, left] {step holds?} (verd);
\draw[refuse] (exam.west) -- node[reflabel, above, align=center] {no path back: no item,\\grade or verdict} (gen.east);
\end{tikzpicture}
\caption{The learner and the independent examiner. The only channel is one-way: per lineage and
generation the learner writes the list of every program the candidate's chain trained on and then the
candidate record (base, adapter and their digests). The examiner holds the exam, drops items that
match a trained program (exactly, or after removing comments and whitespace), serves every model itself
on one pinned build, grades by the direct-completion protocol, applies the registered per-step rule
and signs aggregates only, which are published with the learner's records after the series. The
dotted path does not exist.}\label{fig:examiner}
\end{figure}

\paragraph{What is compared.} Every comparison is within a lineage and between consecutive
generations, on the items no compared model trained on, so a model's gain is measured against the
model it was trained from. The registered decision is per lineage. The four lineages share the exam's
items, so a sign test pooled over lineages treats correlated pairs as independent; the pooled figures
are reported as a description, not as a test. The examiner's base-against-base null replicate, run before RSI4, answered
12 of the 889 items on both sides with 0 discordant items and 0 failed requests: the protocol is
deterministic at the serving build and tier, so a discordant item is a difference between the
models.\cl{RS07}

\subsection{RSI1: self-taught generations, judged first pooled and then per lineage}\label{sec:rsi1}

RSI1 is one registration with dated arms, each added before it ran; its exams were selected and
graded by the arm's own registered analysis script, without a separate examiner.\cl{RS01} Arms A and B
ran the closed-book loop of \cref{sec:documentation} for three generations and were judged on the
loop's own probes; re-graded under a neutral request, most of the gain they measured was adaptation
to the probing protocol and did not transfer.\cl{RS11} Arm H therefore built the exam described
above and measured every generation on it greedily under the neutral request. Its three replicates,
trained only on tasks their own models wrote, answered 38, 31 and 34 of the 889 items at
generation~1, 60, 63 and 49 at generation~2 and 67, 83 and 59 at generation~3 (base 0); pooled over the
replicates the steps were 103/0, 95/26 and 103/66 wins/losses ($p=0.00273$ at the third) with no
replicate falling, which met the registered pooled criterion.\cl{RS12} The same design in Python
(arm HP, 820 items, base 2) met its first two steps (49/7 and 83/8) and not the third (59/50,
$p=0.222$), where one replicate fell from 69 to 60.\cl{RS13} Arm HR added replay across generations
and a pooled retention bound and met the pooled criterion (66/4, 16/1, 25/3), but each step's gain came
from a single replicate, so it is not evidence that one lineage improved twice.\cl{RS14} Three arms
changed the update instead: self-imitation of the model's own verified reasoning (R) and policy
gradient with negative rewards at two update scales (P, P2). After three iterations they answered 21,
4 and 12 of the 889 items, and on six families none of them trained on their pass@1 went from 0.41
to 0.4479, 0.4408 to 0.4242 and 0.4419 to 0.4408: none met its criterion.\cl{RS15}

Because a pooled test cannot tell staggered single gains from one model improving repeatedly, arm HL
repeated arm H's recipe in four fresh lineages under a per-lineage criterion: a lineage succeeds when
all three of its steps hold, and the arm when two of four do. Its first run ended on the project disk's
inode quota before any lineage finished and was not evaluated.\cl{RS16} The rerun met the criterion in
1 of 4 lineages: r3 rose at every step (16/0; 14/5, $p=0.0318$; 33/12, $p=0.00123$), r2 rose twice and
fell at generation~3 (27/36), r0's later steps were not significant (35/25, $p=0.123$; 37/28,
$p=0.161$), and r1 accepted no round after generation~1. Pooled, the third step was 97/76
($p=0.0641$).\cl{RS17} RSI1 thus showed large first steps on items the base never answered and
third steps that did not hold within a lineage.

\subsection{RSI2: on-policy policy gradient}\label{sec:rsi2}

RSI2 asked whether replacing the update with on-policy policy gradient under powered gates makes the
later steps hold. Each generation trained on the current policy's own graded attempts at programs it
solves sometimes ($k=8$ attempts, group-standardised advantages, clip edges 0.2 and 0.28, learning
rate $6\times10^{-6}$ for 24 steps, LoRA rank 16), and the exam moved to the independent
examiner.\cl{RS10,RS18} At generation~1 the four lineages won 4, 2, 3 and 3 items against the base and
lost none ($p$ from 0.0625 to 0.25): no step held, so no lineage could succeed, and the examiner's
plan verdict reads not met.\cl{RS19} The examiner's overlap audit found that the policy's own
proposals had reproduced 1, 1 and 2 exam programs exactly in three lineages; RSI2 had no exclusion
and those items count as scored.\cl{RS19} Generation~2 could not be served, because every trained
candidate carried the same adapter card id; the criterion was already decided, so it was not rerun
after the fix.\cl{RS08} The recipe raised the learner's own dev set in its pilot but moved the exam by
two to four items.

\subsection{RSI3: training on the interpreter's answers to the policy's own programs}\label{sec:rsi3}

RSI3 returned to arm H's source of information, training on the interpreter's output for programs the
policy itself proposed, and replaced what had failed in RSI1: acceptance on a dev set powered for an
effect of 0.08, replay of earlier generations, proposals diversified per manual section, and the whole
learner behind Brain's API. A generation proposes up to 600 programs, predicts each greedily and
trains a LoRA continued from the champion (rank 16, learning rate $3\times10^{-5}$, 40 steps) on the
mispredicted programs, with the correctly predicted ones and earlier generations' records each drawn
at a share of 0.25; the examiner now excludes overlapping items.\cl{RS09,RS05}

Every lineage improved at generations~1 and~2; at generation~3 r2 and r3 improved again, r0 held
flat (29/28, $p=0.5$) and r1 fell (18/28) (\cref{tab:rsi}). Two of four lineages succeeded and the
plan verdict reads not met; pooled, the steps were 288/0, 199/60 and 142/92, and no exam item matched a
training program.\cl{RS21,RS22} The learner's records show the signal running out: over the four
lineages the champions mispredicted 2127 of 2400 proposals at generation~1, 1125 of 2228 at
generation~2 and 860 of 1942 at generation~3, so later generations trained mostly on what their
predecessor already knew, while each also continued the previous adapter.\cl{RS23} The learner's dev
gate accepted all 12 steps, and the examiner found 2 of them not to hold: r0's and r1's third steps,
the two smallest dev margins (13/8 and 7/5).\cl{RS24}

\subsection{RSI4: cumulative retraining from the base on the policy's frontier}\label{sec:rsi4}

RSI4 changed the two mechanisms RSI3's records point at, and with them its training schedule
(\cref{fig:rsi-cycle}).
\emph{Cumulative retraining from the base}: generation $g$'s candidate is a fresh LoRA trained from
the base on the union of every accepted generation's frontier programs, not a continuation of
generation $g-1$'s adapter. This is expert iteration as ReST-EM states it, generating samples,
filtering them by binary feedback and fine-tuning on them, repeated;~\citep{restem2023}\cl{CT16} the data of generation $g$ is still proposed and answered
by generation $g-1$'s champion. \emph{Frontier selection}: the champion predicts each of its proposals
4 times at temperature 1.0, and a program enters training when at least one prediction is wrong, with
the interpreter's output as the target. The frontier therefore holds programs the champion never gets
right and also programs it usually gets right: one it predicts right with probability 0.9 per sample
enters with probability 0.34, one at 0.5 with 0.94.\cl{RS25,RS35} \emph{Schedule}: proposals continue
until the frontier holds 300, 600 and 1200 new programs at generations 1, 2 and 3; training makes 2
passes over the union, so the steps grow with it (the larger of the recipe's steps and
$\lceil 2\times\text{records}/8\rceil$), and there are no separate draws of correctly predicted
programs or replay, the union being the replay.\cl{RS25} The schedule grows geometrically because the
dev pilot's second step fell to a gain of 0.045 on a fixed 300.\cl{RS25} In the run the generations
trained on 300, 900 and 2100 cumulative records for 75, 225 and 525 steps in each lineage.\cl{RS26}
These changes were made together, and no ablation separates their contributions.

\begin{figure}[tbp]
\centering
\begin{tikzpicture}[fig]
\tikzset{cell/.style={text width=28mm}}
\def\dx{40mm}
% row 0: the source
\node[data, cell] (manual) at (\dx,0) {Lua 5.1 manual,\\one section per request};
% row 1
\node[proc=okSky!22, cell] (champ) at (0,-21mm) {champion of\\generation $g-1$\\(the base at $g=1$)};
\node[proc=okSky!22, cell] (prop) at (\dx,-21mm) {champion proposes\\a program (own\\word and number)};
\node[verif, cell] (lua) at (2*\dx,-21mm) {Lua 5.1.5 runs it\\twice; the equal\\output is the label};
\node[proc=okSky!22, cell] (pred) at (3*\dx,-21mm) {champion predicts\\the output $k=4$\\times, $T=1.0$};
% row 2, right to left
\node[decide, aspect=2.2, inner sep=0.5pt, text width=15mm] (front) at (3*\dx,-49mm) {any of the\\4 wrong?};
\node[data, cell] (union) at (2*\dx,-49mm) {frontier union:\\300, 900, 2100\\programs};
\node[proc=okSky!22, cell] (lora) at (\dx,-49mm) {fresh LoRA from\\the base, 2 passes:\\75, 225, 525 steps};
\node[verif, cell] (gate) at (0,-49mm) {learner's gate:\\dev 400 units,\\retention 300};
% row 3
\node[data, cell] (hand) at (\dx,-76mm) {hand-off: every\\trained program,\\candidate record};
% flows
\draw[flow] (champ) -- (prop);
\draw[flow] (manual) -- (prop);
\draw[flow] (prop) -- (lua);
\draw[flow] (lua) -- (pred);
\draw[flow] (pred) -- (front);
\draw[flow] (front) -- node[lab, above] {yes} (union);
\draw[flow] (union) -- (lora);
\draw[flow] (lora) -- (gate);
\draw[flow] ($(gate.south)+(6mm,0)$) |- node[lab, pos=0.3, right] {accepted} (hand.west);
\draw[ctrl] (gate) -- node[lab, right, align=left] {new\\champion} (champ);
\draw[refuse] ($(gate.south)+(-6mm,0)$) -- ++(0,-9mm) node[reflabel, below, align=center] {rejected:\\champion stays};
\draw[refuse] (front.south) -- ++(0,-8mm) node[reflabel, below] {all 4 right: not trained};
\draw[refuse] (lua.north) -- ++(0,6mm) node[reflabel, above, align=center] {refused: no stable output,\\a repeat, a nondeterministic call};
\draw[ctrl] (hand.east) -- node[lab, above] {to the examiner} ++(30mm,0);
\end{tikzpicture}
\caption{One RSI4 generation, read clockwise from the champion. The champion proposes programs about
the manual and the interpreter labels them; a program joins the frontier when the champion predicts it
wrong in at least one of 4 sampled attempts (300, 600 and 1200 new programs at generations 1 to 3); the
candidate is a fresh LoRA trained from the base on the union of every generation's frontier; the
learner's gate decides whether it becomes the next champion, and the examiner, not the gate, decides
the registered outcome. Dotted arrows are what each stage refuses.}\label{fig:rsi-cycle}
\end{figure}

\paragraph{Before the run.} Three things were settled before any verdict existed and are recorded
with the launch. The null replicate above confirmed the protocol. The base was served at the f32 KV
tier, at which it answers 12 items; RSI3's stored baseline had been served at the build's default tier,
at which it answered none, so RSI4's levels are not comparable with RSI3's, only its steps
are.\cl{RS07,RS27} And the examiner's build was changed before the first verdict because its first
build would have marked every verdict invalid: the learner's Brain declares the base digest over the
weight files alone and the examiner's over configuration, weights and index, both computed from the
same files; the new build matches a declared base by the checkpoint's content under either rule.
Exam, protocol, serving build, tier and criterion were unchanged.\cl{RS06}

\begin{table}[tbp]
\centering\small
\setlength{\tabcolsep}{3.4pt}
\input{generated/rsi_table}
\caption{RSI3 and RSI4 as the examiner decided them: items right of 889 at the base and each
generation, and each step's wins/losses against its predecessor with the exact one-sided sign test
over the items no compared model trained on (in RSI4's r2 one item was excluded). A step holds at
$p<0.05$ with wins above losses; the series' criterion is that at least three lineages hold at every
step. Pooled rows sum the lineage rows (they equal the plan verdict's) and are descriptive, without a test, because the lineages share the exam's items. RSI3's base was
served at a different KV tier from RSI4's, so the levels of the two series differ by setup; compare
steps, not levels.}\label{tab:rsi}
\end{table}

\begin{figure}[tbp]
\centering
\begin{tikzpicture}
\begin{groupplot}[paper, group style={group size=2 by 1, horizontal sep=12mm},
  width=0.48\linewidth, height=5.4cm, xmin=-0.25, xmax=3.25, ymin=0, ymax=260,
  xtick={0,1,2,3}, xticklabels={base,1,2,3}, xlabel={generation},
  ylabel={items right of 889 (exam subset)}]
\nextgroupplot[title={RSI3: continued adapter (not met)}]
\addplot[lineage r0] table[col sep=comma,x=gen,y=r0] {generated/rsi3_items.csv};
\addplot[lineage r1] table[col sep=comma,x=gen,y=r1] {generated/rsi3_items.csv};
\addplot[lineage r2] table[col sep=comma,x=gen,y=r2] {generated/rsi3_items.csv};
\addplot[lineage r3] table[col sep=comma,x=gen,y=r3] {generated/rsi3_items.csv};
\nextgroupplot[title={RSI4: retrained from the base (met)}, ylabel={},
  legend style={at={(-0.11,-0.3)}, anchor=north, legend columns=4, /tikz/every even column/.append style={column sep=4mm}}]
\addplot[lineage r0] table[col sep=comma,x=gen,y=r0] {generated/rsi4_items.csv};
\addplot[lineage r1] table[col sep=comma,x=gen,y=r1] {generated/rsi4_items.csv};
\addplot[lineage r2] table[col sep=comma,x=gen,y=r2] {generated/rsi4_items.csv};
\addplot[lineage r3] table[col sep=comma,x=gen,y=r3] {generated/rsi4_items.csv};
\legend{lineage r0, lineage r1, lineage r2, lineage r3}
\end{groupplot}
\end{tikzpicture}
\caption{Items right on the sealed exam per generation and lineage, as the examiner measured them; the
exam holds only items the base could not answer, so these are counts on that hard subset, not an
accuracy on Lua programs. In RSI3 two lineages stop rising at generation~3 (r0 flat, r1 down); in RSI4
every lineage rises at every generation. The base level differs between the panels because the base
was served at a different KV tier (\cref{tab:rsi}).}\label{fig:rsi-items}
\end{figure}

\begin{figure}[tbp]
\centering
\begin{tikzpicture}
\begin{groupplot}[paper, group style={group size=2 by 1, horizontal sep=12mm},
  width=0.48\linewidth, height=5.2cm, ybar=0pt, /pgf/bar width=3.4pt, /pgf/bar shift=0pt,
  ymin=-70, ymax=140, xmin=-1, xmax=15,
  xtick={1,5,9,13}, xticklabels={r0,r1,r2,r3}, xlabel={lineage (steps 1, 2, 3 left to right)},
  ylabel={items won ($+$) and lost ($-$)}, extra y ticks={0},
  extra y tick labels={}, extra y tick style={grid style={black!70, thin}}]
\nextgroupplot[title={RSI3},
  legend style={at={(1.06,-0.3)}, anchor=north, legend columns=3, /tikz/every even column/.append style={column sep=4mm}}]
\addplot[fill=okBlue, draw=okBlue!70!black, area legend] table[col sep=comma,x=x,y=wins] {generated/rsi3_steps.csv};
\addplot[area legend, fill=white, draw=okVermillion!80!black, postaction={pattern=north east lines, pattern color=okVermillion}] table[col sep=comma,x=x,y=losses] {generated/rsi3_steps.csv};
\addplot[only marks, mark=x, mark size=3.2pt, very thick, black] table[col sep=comma,x=x,
  y expr={\thisrow{holds} < 0.5 ? \thisrow{wins} + 14 : nan}] {generated/rsi3_steps.csv};
\legend{items won, items lost, step does not hold}
\nextgroupplot[title={RSI4}, ylabel={}]
\addplot[fill=okBlue, draw=okBlue!70!black, area legend] table[col sep=comma,x=x,y=wins] {generated/rsi4_steps.csv};
\addplot[area legend, fill=white, draw=okVermillion!80!black, postaction={pattern=north east lines, pattern color=okVermillion}] table[col sep=comma,x=x,y=losses] {generated/rsi4_steps.csv};
\addplot[only marks, mark=x, mark size=3.2pt, very thick, black] table[col sep=comma,x=x,
  y expr={\thisrow{holds} < 0.5 ? \thisrow{wins} + 14 : nan}] {generated/rsi4_steps.csv};
\end{groupplot}
\end{tikzpicture}
\caption{Wins and losses of every step against its predecessor (three bars per lineage: generation 1
against the base, 2 against 1, 3 against 2); a cross marks a step that does not hold. Losses grow with
the generation in both series, because a better model has more items to lose; in RSI4 wins outgrow
them at every step, in RSI3 they do not at r0's and r1's third steps.}\label{fig:rsi-steps}
\end{figure}


\paragraph{Outcome: the criterion is met in all four lineages.} Every one of the twelve RSI4 steps
held (\cref{tab:rsi,fig:rsi-items,fig:rsi-steps}). From 12 items at the base the lineages reached 92,
131, 118 and 121 at generation~1, 138, 181, 169 and 187 at generation~2 and 192, 225, 241 and 218 at
generation~3; the weakest step, r3's third, won 67 items and lost 36 ($p=0.00146$).\cl{RS28} The plan
verdict reads met, with 4 of 4 lineages succeeding where the rule needed 3.\cl{RS29} Summed over the
lineages, as a description, the steps won and lost 422/8, 320/107 and 354/153 items.\cl{RS30} One exam item matched a training program of lineage r2 and was
excluded from r2's comparisons; no other overlap and no integrity flag was raised.\cl{RS31} No lineage
stalled at generation~3. Each lineage's learner
took 2.4 to 3.4 hours of one card for all three generations.\cl{RS26}

\paragraph{The learner's gate did not predict the exam.} RSI4's dev gate accepted all 12 steps, and
all 12 held on the exam; but its evidence was weakest where the exam's was not. At r1's third step
the dev set showed 25 wins and 19 losses ($p=0.226$) while the exam showed 99 and 55
($p=0.00024$).\cl{RS33} Together with RSI3, where the same greedy gate accepted two steps the exam
rejected, this says that a 400-unit dev set reused across generations is too weak an instrument to
decide a generation on. TRUST replaces such a greedy gate with an online multiple-testing procedure that controls the rate
of accepted inferior updates across rounds;~\citep{trust2026}\cl{CT15} a gate of that kind now exists in the learner as an
option; replayed on RSI3's dev counts at a tolerance of 0.02 it would have kept the incumbent at both
of RSI3's failed steps most of the time, at the cost of rejecting some true gains. No RSI run used
it.\cl{RS34}

\subsection{What the series establishes}\label{sec:rsi-scope}

On one task family, predicting the output of Lua~5.1 programs, with one base model, Qwen3-4B, and no
stronger model, teacher or human answer in training, four independent lineages of self-taught
generations, each retrained from the base on the interpreter-verified frontier of its own proposals,
improved at each of three consecutive generations on an exam the learner never read, as judged by an
independent examiner under a criterion fixed before the run. The gain is measured on items the base
could not answer, the labels come from the interpreter, and each lineage is judged on its own. The
earlier series locate what that needed: answers from the environment rather than from the policy's rewards (RSI1's arms R, P and P2,
RSI2), a per-lineage test rather than a pooled one (arms HR and HL), and, at the third generation,
retraining from the base on the policy's frontier rather than continuing its adapter on whatever it
proposed (RSI3 against RSI4). The series does not establish transfer to another task family or
model, behaviour beyond three generations, which of RSI4's simultaneous changes (retraining from the
base, frontier selection, the growing schedule) carries the gain, or that the improvement compounds
because an improved model proposes better material: that needs the same frontier data proposed by the
base or a fixed policy as a control, which was not run.
